\documentclass[10pt,twocolumn]{article}
\usepackage[a4paper,margin=1.9cm,columnsep=0.7cm]{geometry}
\usepackage{mathptmx}
\usepackage[T1]{fontenc}
\usepackage{microtype}
\usepackage{graphicx}
\usepackage{booktabs}
\usepackage{caption}
\usepackage{authblk}
\usepackage[hidelinks]{hyperref}
\usepackage{xcolor}

\captionsetup{font=small,labelfont=bf,skip=5pt}
\setlength{\parindent}{1.1em}
\definecolor{mono}{HTML}{2F6F9F}
\definecolor{dia}{HTML}{C25A3C}

\title{\vspace{-1.2em}\bfseries Monologue or Dialogue?\\[2pt]
\large Measuring Conversational Structure in the Reasoning Traces of Large Language Models}
\author[ ]{Oskar Engl}
\affil[ ]{\small Johannes Kepler University Linz, Austria \quad\textit{oskar.engl@ik.me}}
\date{\small Condensed version --- July 2026}

\begin{document}
\twocolumn[
\begin{@twocolumnfalse}
\maketitle
\begin{abstract}
\noindent
A reasoning model such as DeepSeek-R1 emits a long ``thinking'' trace before it answers. That
trace can read as one mind reasoning step by step, or as several voices trading proposals and
objections. We ask which description the \emph{text} fits when the question is put to measurement.
Because no reasoning trace comes labelled with its true voice structure, we build a $2\times2$
grid of corpora whose structure \emph{is} known --- human and AI, monologue and dialogue --- and
let each detection method judge the traces only after it passes a pre-registered validity gate on
the four known corpora. Of three methods, two style-based detectors pass and a dialogue-act
detector fails and is set aside. The two survivors \emph{disagree}: a document-level author-count
method places the traces just onto the monologue side of its (wide) decision band, while a
sentence-level style-switch method places them firmly past the dialogue pole. R1 traces therefore
carry local, sentence-level marks of style switching \emph{without} a matching separation into
distinct document-level authors --- a single narrator whose register shifts, not a transcribed
conversation.
\medskip
\end{abstract}
\end{@twocolumnfalse}
]

\begin{figure*}[t]
\centering
\includegraphics[width=0.94\textwidth]{fig_grid.pdf}
\caption{\textbf{The design.} Left to right: four corpora of known voice structure feed a
\emph{validity gate} that each of the three detection methods must pass --- separating monologue
from dialogue on all four known cells --- before it is allowed to judge the target. Only surviving
methods are applied to the DeepSeek-R1 traces (\texttt{ai-reason}) and place them on the
monologue--dialogue axis anchored by the two same-task AI cells; the human cells validate only.}
\label{fig:grid}
\end{figure*}

\section{The question}
Reasoning-oriented training is now standard in frontier models: given a problem, a model like
DeepSeek-R1 \cite{guo2025} first generates an extended stretch of intermediate tokens
\cite{kambhampati2025} --- a reasoning trace --- and only then commits to an answer. These traces
have a curious property. Each is produced by a single model in one pass, yet a trace often contains
questions and their answers, objections and concessions, and abrupt reversals
(``Wait, that's not right''). In human writing such features can mark a genuine exchange, but they
can equally mark one vivid inner voice. Hence a question easy to state and less easy to settle:
\emph{does a reasoning trace have the textual structure of a monologue or of a dialogue?}

The question has a live theoretical stake. Kim et al.~\cite{kim2026} argue that reasoning-trained
models improve partly because they simulate a ``society of thought'' --- distinct internal
perspectives that propose, criticise, and reconcile --- and older ideas echo it: reasoning as at
root argumentative and social \cite{mercier2017}, and the mind as many interacting agents
\cite{minsky1986}. Our premise is narrower and testable: voice structure leaves marks in the
writing itself. If distinct perspectives take turns producing the text, they should also differ in
the small, largely involuntary habits --- function-word use, sentence rhythm, dialogue-act moves
--- that authorship analysis uses to tell writers apart. We therefore treat the question
operationally: a monologue should read as one author, a dialogue as several, and we ask which the
traces resemble.

\section{The design: a calibration grid with a gate}
No corpus of reasoning traces is labelled with its true voice structure, which is what makes the
question hard to settle directly. We address this with the grid in Figure~\ref{fig:grid}. Four
corpora of \emph{known} structure vary producer (human vs.\ AI) and structure (monologue vs.\
dialogue): human monologue (think-aloud protocols \cite{wurgaft2025}), human dialogue (group
deliberation, DeliData \cite{karadzhov2023}), AI monologue (one Qwen-2.5-7B instance solving a
puzzle alone), and AI dialogue (four stateless Qwen-2.5-7B instances, each given a fixed role,
debating). The fifth cell --- native DeepSeek-R1 traces --- is the target.

The grid works as a \emph{gate}. A method may report on the target only after it correctly
separates monologue from dialogue on all four known cells; a method that fails is reported as
failing and set aside. A detector that cannot tell voices apart in a genuine four-way debate is in
no position to rule them out in a trace; one that reports several voices in a lone think-aloud
protocol is in no position to affirm them there.

Two design choices keep the contrast clean. First, \emph{task control}: \texttt{ai-mono},
\texttt{ai-dia}, \texttt{ai-reason} and \texttt{hum-mono} all use the same puzzle, Game of 24
\cite{yao2023}, so a method reading a monologue as one voice cannot merely be reacting to subject
matter. (\texttt{hum-dia} uses the Wason task, the closest human group-deliberation data available;
its role is only to confirm that a method detects human dialogue at all, so it never anchors the
placement.) Second, \emph{pre-registration}: every method configuration, gate, threshold and figure
was fixed in a written specification and committed to version control (commit \texttt{8cbdee4},
2026-07-17) before any statistic was computed on the traces. Speaker-identity tokens
(pseudonyms, role titles, letter labels) were stripped from both dialogue cells beforehand, so
dialogue detection cannot be an artefact of speaker cues.

\section{Three methods, three signals}
The methods read deliberately different signals, so that agreement among them would carry more
weight than any one alone.

\textbf{Method 1 --- document-level author count} \cite{nath2019}. Clusters a document's
paragraph-sized windows on their 50-most-frequent-token profiles and reports the number of
clusters as an author count, reduced to a binary single/multi call. \emph{Gate:} balanced accuracy
across the four cells must beat the 0.50 label-permutation baseline.

\textbf{Method 2 --- sentence-level style switch} \cite{schmidt2025}. A pretrained sequential
sentence-pair classifier emits, for each adjacent sentence pair, the probability that the two were
written by different authors, using content-independent style embeddings \cite{patel2025}. Because
raw scores depend on length, each document's score is turned into a length-adjusted \emph{lean}
--- its distance to the dialogue anchor's fitted length curve minus its distance to the monologue
anchor's. \emph{Gate:} mean lean must be higher on the dialogue cells than the monologue cells,
with Cliff's $\delta\ge0.33$.

\textbf{Method 3 --- dialogue-act structure} \cite{duran2022,schegloff1973}. An LLM tags each
utterance with a CAMS dialogue act and measures adjacency-pair completion: when a question,
request or offer occurs, is it met by a valid answer or acceptance in the very next utterance ---
the mark of conversation a lone writer has little reason to produce. \emph{Gates:} the completion
rate must separate dialogue from monologue cells \emph{and} the tagger's labels must reach
Cohen's $\kappa\ge0.70$ against hand labels.

\section{Results: a split}
\textbf{Method 1 passes, but softly.} Balanced accuracy is $0.6088$ (bootstrap 95\% CI
$[0.596,0.621]$; permutation $p=0.0001$), clearing the baseline. But the pass leans on
near-perfect single-author recall ($0.975$ on \texttt{hum-mono}, $0.998$ on \texttt{ai-mono})
averaged against weak multi-author recall: $0.423$ on \texttt{ai-dia} and a complete miss on
\texttt{hum-dia} ($0.000$). The detector flags dialogue where it finds it but essentially never
rules single authorship \emph{out}. A monologue-leaning placement under this method should thus be
read as the \emph{absence} of a detected authorship signal, not positive evidence of one voice.

\textbf{Method 2 passes cleanly.} Mean lean is higher on the dialogue cells than the monologue cells
($0.0125$ vs.\ $-0.0105$) with Cliff's $\delta=0.6211$, well above the $0.33$ floor. The traces
themselves ($0.0433$) sit above \emph{both} AI anchors.

\textbf{Method 3 fails its reliability gate.} Its completion rate does separate dialogue from
monologue (Cliff's $\delta=0.4717$), but the tagger reaches only $\kappa=0.5425$
(95\% CI $[0.479,0.604]$) against hand labels --- below the pre-registered $0.70$, with the interval
not even touching it. Much of the disagreement is between semantically adjacent tags
(a tree-weighted $\alpha$ reaches $0.619$), but under the frozen rule the method is reported as a
failed validation and takes no part in placing the target.

\begin{figure*}[t]
\centering
\includegraphics[width=0.90\textwidth]{fig_split.pdf}
\caption{\textbf{The answer in one view.} Each surviving method's placement of the target on the
monologue--dialogue axis (poles at $0$ and $1$). The shaded band is that method's margin, within
which no lean is called. Method~1's marker sits at the edge of a \emph{wide} margin, just onto the
monologue side --- a marginal call the bootstrap reverses in nearly a third of resamples.
Method~2 lands well past the dialogue pole (stable in every resample). The two calls disagree ---
the result is a split.}
\label{fig:split}
\end{figure*}

\textbf{The two survivors disagree} (Table~\ref{tab:verdict}, Figure~\ref{fig:split}). Under the
placement rule fixed in advance, Method~1 puts the traces just onto the \emph{monologue} side of
its own no-decision band ($r=0.410$ vs.\ margin $m=0.361$) --- a marginal call the bootstrap
reverses in nearly a third of resamples. Method~2 puts them firmly on the \emph{dialogue} side,
past the dialogue anchor itself ($r=0.737$ vs.\ $m=0.112$; stable in every resample). By the rule
fixed in advance this is reported as a split, not resolved by a vote.

\begin{table}[t]
\centering
\small
\caption{Cross-method verdicts on the reasoning traces. This is the answer to the research
question.}
\label{tab:verdict}
\begin{tabular}{@{}lcc@{}}
\toprule
Method & Passed gate? & Verdict on \texttt{ai-reason} \\
\midrule
1: author count (doc.) & pass & \textcolor{mono}{monologue-leaning} \\
2: style switch (sent.) & pass & \textcolor{dia}{dialogue-leaning} \\
3: dialogue acts & \textbf{fail} & excluded --- no reading \\
\bottomrule
\end{tabular}
\end{table}

\section{What the split means}
The claim this study licenses is about \emph{textual surface structure}, not hidden minds: whether
the traces exhibit the style-switching and adjacency-pair regularities of multi-party dialogue.
Read together, the two survivors describe the traces as text that carries local, sentence-level
marks of style switching \emph{without} a matching separation into distinct document-level authors.
Two qualifications both pull toward the middle: Method~1 is a conservative detector (it recovers
multi-voice structure in under a quarter of genuinely multi-voice text), so its monologue call is
weak evidence; and Method~2's monologue distance is \emph{extrapolated}, because the traces run far
longer than the monologue anchor --- what it measures directly is that the traces track the
\emph{dialogue} anchor closely across the shared length range.

The upshot is that the outside view and the inside view do not straightforwardly coincide. Kim et
al.~\cite{kim2026} infer distinct internal perspectives from interpretability and trace labelling;
the surface text measured here does not resolve into distinct authorial voices even where its
register shifts. Internal multi-agent dynamics can coexist with a stylometrically single-author
surface, and the present evidence is consistent with that gap --- not with a text that reads, on its
face, as a transcribed conversation. We report this as the reading the pre-registered frame
supports, not as a refutation of the internal account.

\section{Limitations}
The target is one model (DeepSeek-R1) on one puzzle family, so the verdict is about this cell, not
reasoning traces at large. The target is also a \emph{different} base model than the anchors
(Qwen-2.5-7B), so any measured distance carries a shared model offset; the placement is best read
as ``closer to which pole,'' not as an absolute distance. Method~2's monologue comparison rests on
extrapolation beyond the fitted length range, and all three methods run out of domain. These are the
reasons the validity gate exists: a method that survives it has at least shown competence on this
data.

\begin{thebibliography}{99}
\small
\bibitem{guo2025} D.\ Guo et al. DeepSeek-R1 incentivizes reasoning in LLMs through reinforcement
learning. \textit{Nature} 645:633--638, 2025.
\bibitem{kambhampati2025} S.\ Kambhampati et al. Stop anthropomorphizing intermediate tokens as
reasoning/thinking traces! \textit{arXiv:2504.09762}, 2025.
\bibitem{kim2026} J.\ Kim, S.\ Lai, N.\ Scherrer, B.\ Ag\"uera y Arcas, J.\ Evans. Reasoning models
generate societies of thought. \textit{arXiv:2601.10825}, 2026.
\bibitem{mercier2017} H.\ Mercier, D.\ Sperber. \textit{The Enigma of Reason.} Harvard Univ.\
Press, 2017.
\bibitem{minsky1986} M.\ Minsky. \textit{The Society of Mind.} Simon \& Schuster, 1986.
\bibitem{wurgaft2025} D.\ Wurgaft et al. Scaling up the think-aloud method. \textit{Proc.\ CogSci},
2025.
\bibitem{karadzhov2023} G.\ Karadzhov, T.\ Stafford, A.\ Vlachos. DeliData: a dataset for
deliberation in multi-party problem solving. \textit{Proc.\ ACM HCI} 7(CSCW2), 2023.
\bibitem{yao2023} S.\ Yao et al. Tree of thoughts: deliberate problem solving with LLMs.
\textit{NeurIPS} 36, 2023.
\bibitem{nath2019} S.\ Nath. Style change detection by threshold-based and window-merge clustering.
\textit{CLEF Working Notes}, 2019.
\bibitem{schmidt2025} G.\ Schmidt et al. Style change detection using a sequential sentence-pair
classifier. \textit{CLEF Working Notes}, 2025. \textit{arXiv:2508.00675}.
\bibitem{patel2025} A.\ Patel et al. StyleDistance: content-independent style embeddings with
synthetic parallel examples. \textit{NAACL}, 2025.
\bibitem{duran2022} N.\ Duran, S.\ Battle, J.\ Smith. Inter-annotator agreement using the
Conversation Analysis Modelling Schema. \textit{Comm.\ Methods \& Measures} 16(3):182--214, 2022.
\bibitem{schegloff1973} E.\ Schegloff, H.\ Sacks. Opening up closings. \textit{Semiotica}
8(4):289--327, 1973.
\end{thebibliography}

\end{document}
